\documentclass[11pt,a4paper]{article}

\usepackage[utf8]{inputenc}
\usepackage[T1]{fontenc}
\usepackage[french]{babel}
\usepackage[a4paper,margin=2.2cm]{geometry}
\usepackage{amsmath,amssymb,amsthm,mathtools}
\usepackage{lmodern}
\usepackage{enumitem}
\usepackage{hyperref}
\usepackage{xcolor}
\usepackage{fancyhdr}
\usepackage{stmaryrd}
\usepackage{mathrsfs}
\usepackage{array}

\hypersetup{
    colorlinks=true,
    linkcolor=black,
    urlcolor=blue,
    citecolor=black
}

\setlength{\parindent}{0pt}
\setlength{\parskip}{0.6em}

\pagestyle{fancy}
\fancyhf{}
\fancyhead[L]{Algèbre linéaire --- Chapitre 11}
\fancyhead[R]{Réduction d’endomorphismes}
\fancyfoot[C]{\thepage}

% Environnements
\newtheorem{definition}{Définition}[section]
\newtheorem{theoreme}[definition]{Théorème}
\newtheorem{proposition}[definition]{Proposition}
\newtheorem{corollaire}[definition]{Corollaire}
\newtheorem{lemme}[definition]{Lemme}
\theoremstyle{remark}
\newtheorem{remarque}[definition]{Remarque}
\newtheorem{exemple}[definition]{Exemple}
\newtheorem*{methode}{Méthode}
\newtheorem*{idee}{Idée}

% Commandes
\newcommand{\R}{\mathbb{R}}
\newcommand{\C}{\mathbb{C}}
\newcommand{\K}{\mathbb{K}}
\newcommand{\N}{\mathbb{N}}
\newcommand{\Z}{\mathbb{Z}}
\newcommand{\Vect}{\mathrm{Vect}}
\newcommand{\Ker}{\mathrm{Ker}}
\newcommand{\Img}{\mathrm{Im}}
\newcommand{\rg}{\mathrm{rg}}
\newcommand{\id}{\mathrm{Id}}
\newcommand{\tr}{\mathrm{tr}}
\newcommand{\Sp}{\mathrm{Sp}}
\newcommand{\diag}{\mathrm{diag}}
\newcommand{\Mnn}{M_n(\K)}
\newcommand{\GLn}{GL_n(\K)}
\newcommand{\chiA}{\chi_A}

\begin{document}

\begin{center}
    {\LARGE \textbf{Chapitre 11 --- Réduction d’endomorphismes}}\\[0.4em]
    {\large Valeurs propres, diagonalisation et trigonalisabilité}
\end{center}

\hrule
\vspace{1em}

\tableofcontents

\newpage

\section{Pourquoi réduire un endomorphisme ?}

Lorsqu’on étudie un endomorphisme
\[
u:E\to E
\]
ou une matrice carrée \(A\), on cherche souvent à le comprendre le plus simplement possible.

L’idée de la \textbf{réduction} est la suivante :
\begin{center}
\textit{peut-on choisir une base dans laquelle la matrice de \(u\) devient très simple ?}
\end{center}

Les formes les plus simples sont :
\begin{itemize}
    \item la forme diagonale ;
    \item la forme triangulaire ;
    \item plus tard, selon le programme, des formes plus fines.
\end{itemize}

\begin{idee}
Réduire un endomorphisme, c’est trouver un bon point de vue.
\end{idee}

\begin{remarque}
La diagonalisation est idéale, car elle simplifie énormément les calculs :
\[
A^n,\quad P(A),\quad e^A,\quad \text{etc.}
\]
Mais toutes les matrices ne sont pas diagonalisables.
\end{remarque}

\section{Valeurs propres et vecteurs propres}

\subsection{Définitions}

\begin{definition}
Soit \(E\) un espace vectoriel sur \(\K\), et
\[
u\in\Lin(E,E)
\]
un endomorphisme.

Un scalaire \(\lambda\in\K\) est appelé une \textbf{valeur propre} de \(u\) s’il existe un vecteur non nul \(x\in E\) tel que
\[
u(x)=\lambda x.
\]
Un tel vecteur \(x\) est appelé un \textbf{vecteur propre} de \(u\) associé à la valeur propre \(\lambda\).
\end{definition}

\begin{remarque}
Le vecteur nul n’est jamais considéré comme vecteur propre.
\end{remarque}

\begin{exemple}
L’endomorphisme
\[
u:\R^2\to\R^2,\qquad u(x,y)=(2x,3y)
\]
admet :
\begin{itemize}
    \item la valeur propre \(2\), avec comme vecteurs propres les vecteurs non nuls de la forme \((x,0)\) ;
    \item la valeur propre \(3\), avec comme vecteurs propres les vecteurs non nuls de la forme \((0,y)\).
\end{itemize}
\end{exemple}

\begin{exemple}
Pour l’identité \(\id_E\), tout vecteur non nul est vecteur propre associé à la valeur propre \(1\).
\end{exemple}

\begin{exemple}
Pour l’homothétie \(u=\lambda \id_E\), tout vecteur non nul est vecteur propre associé à \(\lambda\).
\end{exemple}

\subsection{Sous-espace propre}

\begin{definition}
Soit \(\lambda\in\K\). On appelle \textbf{sous-espace propre} de \(u\) associé à \(\lambda\) l’ensemble
\[
E_\lambda(u)=\Ker(u-\lambda\id_E).
\]
Autrement dit,
\[
E_\lambda(u)=\{x\in E\mid u(x)=\lambda x\}.
\]
\end{definition}

\begin{proposition}
\(E_\lambda(u)\) est un sous-espace vectoriel de \(E\).
\end{proposition}

\begin{proof}
Comme
\[
E_\lambda(u)=\Ker(u-\lambda\id_E),
\]
c’est le noyau d’une application linéaire, donc un sous-espace vectoriel.
\end{proof}

\begin{proposition}
\(\lambda\) est une valeur propre de \(u\) si et seulement si
\[
E_\lambda(u)\neq \{0\}.
\]
\end{proposition}

\begin{proof}
Par définition, \(\lambda\) est valeur propre s’il existe un vecteur propre non nul \(x\), ce qui équivaut à dire qu’il existe un élément non nul dans \(E_\lambda(u)\).
\end{proof}

\begin{remarque}
Les vecteurs propres associés à \(\lambda\), auxquels on ajoute \(0\), forment exactement le sous-espace propre \(E_\lambda(u)\).
\end{remarque}

\section{Premières propriétés des valeurs propres}

\subsection{Valeurs propres d’un endomorphisme simple}

\begin{proposition}
Les valeurs propres de l’homothétie \(\lambda\id_E\) sont réduites à \(\lambda\).
\end{proposition}

\begin{proof}
Si
\[
(\lambda \id_E)(x)=\mu x
\]
avec \(x\neq 0\), alors
\[
\lambda x=\mu x
\Longrightarrow
(\lambda-\mu)x=0.
\]
Comme \(x\neq 0\), on obtient \(\mu=\lambda\).
\end{proof}

\subsection{Valeurs propres de matrices triangulaires}

\begin{proposition}
Les valeurs propres d’une matrice triangulaire sont ses coefficients diagonaux.
\end{proposition}

\begin{proof}
Le polynôme caractéristique d’une matrice triangulaire
\[
A=(a_{ij})
\]
vérifie
\[
\chiA(X)=\det(XI_n-A),
\]
et \(XI_n-A\) est encore triangulaire, de diagonale
\[
X-a_{11},\dots,X-a_{nn}.
\]
Donc
\[
\chiA(X)=(X-a_{11})\cdots(X-a_{nn}).
\]
Les racines sont donc exactement les coefficients diagonaux.
\end{proof}

\section{Familles de vecteurs propres}

\subsection{Lemme fondamental}

\begin{theoreme}
Des vecteurs propres associés à des valeurs propres deux à deux distinctes sont linéairement indépendants.
\end{theoreme}

\begin{proof}
Soient \(\lambda_1,\dots,\lambda_p\) des valeurs propres deux à deux distinctes, et \(x_1,\dots,x_p\) des vecteurs propres associés.

Nous allons montrer que la famille \((x_1,\dots,x_p)\) est libre par récurrence sur \(p\).

\textbf{Initialisation :} pour \(p=1\), la famille \((x_1)\) est libre car \(x_1\neq 0\).

\textbf{Hérédité :} supposons le résultat vrai au rang \(p-1\). Supposons
\[
\alpha_1x_1+\cdots+\alpha_px_p=0.
\]
Appliquons \(u\) :
\[
\alpha_1\lambda_1x_1+\cdots+\alpha_p\lambda_px_p=0.
\]
Multiplions la première relation par \(\lambda_p\) et soustrayons :
\[
\alpha_1(\lambda_1-\lambda_p)x_1+\cdots+\alpha_{p-1}(\lambda_{p-1}-\lambda_p)x_{p-1}=0.
\]
Par hypothèse de récurrence, la famille \((x_1,\dots,x_{p-1})\) est libre, donc
\[
\alpha_i(\lambda_i-\lambda_p)=0
\qquad \text{pour } i=1,\dots,p-1.
\]
Comme les \(\lambda_i\) sont distinctes, on a
\[
\lambda_i-\lambda_p\neq 0,
\]
donc
\[
\alpha_1=\cdots=\alpha_{p-1}=0.
\]
La relation initiale devient
\[
\alpha_px_p=0,
\]
et comme \(x_p\neq 0\), on obtient \(\alpha_p=0\). La famille est donc libre.
\end{proof}

\begin{corollaire}
Des sous-espaces propres associés à des valeurs propres deux à deux distinctes sont en somme directe.
\end{corollaire}

\begin{proof}
Si l’on prend une famille de vecteurs, chacun dans un sous-espace propre différent, alors elle est libre par le théorème précédent.
\end{proof}

\begin{corollaire}
Si un endomorphisme de dimension \(n\) possède \(n\) valeurs propres distinctes, alors il est diagonalisable.
\end{corollaire}

\begin{proof}
On choisit un vecteur propre non nul dans chacun des \(n\) sous-espaces propres. On obtient \(n\) vecteurs libres dans un espace de dimension \(n\), donc une base de vecteurs propres.
\end{proof}

\section{Polynôme caractéristique}

\subsection{Définition}

\begin{definition}
Soit \(A\in M_n(\K)\).  
On appelle \textbf{polynôme caractéristique} de \(A\) le polynôme
\[
\chiA(X)=\det(XI_n-A).
\]
\end{definition}

\begin{remarque}
C’est un polynôme de degré \(n\), unitaire.
\end{remarque}

\begin{proposition}
Un scalaire \(\lambda\in\K\) est une valeur propre de \(A\) si et seulement si
\[
\chiA(\lambda)=0.
\]
\end{proposition}

\begin{proof}
\(\lambda\) est valeur propre si et seulement si il existe \(x\neq 0\) tel que
\[
Ax=\lambda x,
\]
c’est-à-dire
\[
(A-\lambda I_n)x=0.
\]
Cela équivaut à dire que la matrice \(A-\lambda I_n\) n’est pas inversible, donc
\[
\det(A-\lambda I_n)=0.
\]
Comme
\[
\det(\lambda I_n-A)=(-1)^n\det(A-\lambda I_n),
\]
cela revient à
\[
\chiA(\lambda)=0.
\]
\end{proof}

\begin{corollaire}
Les valeurs propres d’une matrice sont les racines de son polynôme caractéristique.
\end{corollaire}

\subsection{Exemples}

\begin{exemple}
Soit
\[
A=
\begin{pmatrix}
2 & 0\\
0 & 3
\end{pmatrix}.
\]
Alors
\[
\chiA(X)=\det\begin{pmatrix}
X-2 & 0\\
0 & X-3
\end{pmatrix}
=(X-2)(X-3).
\]
Les valeurs propres sont donc \(2\) et \(3\).
\end{exemple}

\begin{exemple}
Soit
\[
A=
\begin{pmatrix}
0 & -1\\
1 & 0
\end{pmatrix}.
\]
Alors
\[
\chiA(X)=\det\begin{pmatrix}
X & 1\\
-1 & X
\end{pmatrix}
=X^2+1.
\]
Cette matrice n’a pas de valeur propre réelle, mais elle a pour valeurs propres complexes \(i\) et \(-i\).
\end{exemple}

\section{Multiplicités algébrique et géométrique}

\subsection{Définitions}

\begin{definition}
Soit \(\lambda\) une valeur propre de \(A\).

Sa \textbf{multiplicité algébrique} est sa multiplicité comme racine du polynôme caractéristique \(\chiA\).

Sa \textbf{multiplicité géométrique} est
\[
\dim(E_\lambda),
\]
c’est-à-dire la dimension du sous-espace propre associé.
\end{definition}

\subsection{Inégalité fondamentale}

\begin{theoreme}
Pour toute valeur propre \(\lambda\),
\[
1\le \dim(E_\lambda)\le m_\lambda,
\]
où \(m_\lambda\) désigne sa multiplicité algébrique.
\end{theoreme}

\begin{proof}
Le fait que \(\lambda\) soit valeur propre implique \(E_\lambda\neq \{0\}\), donc
\[
\dim(E_\lambda)\ge 1.
\]

Pour montrer
\[
\dim(E_\lambda)\le m_\lambda,
\]
on choisit une base adaptée à \(E_\lambda\), de sorte que la matrice de \(u\) prenne une forme par blocs dans laquelle \(\lambda\) apparaît au moins \(\dim(E_\lambda)\) fois sur la diagonale. Le polynôme caractéristique contient alors un facteur
\[
(X-\lambda)^{\dim(E_\lambda)}.
\]
La multiplicité algébrique est donc au moins égale à la multiplicité géométrique.
\end{proof}

\begin{remarque}
Cette inégalité est fondamentale : la dimension du sous-espace propre ne peut jamais dépasser la multiplicité algébrique.
\end{remarque}

\section{Matrices diagonales et diagonalisation}

\subsection{Matrices diagonales}

\begin{definition}
Une matrice diagonale est une matrice de la forme
\[
D=
\begin{pmatrix}
\lambda_1 & 0 & \cdots & 0\\
0 & \lambda_2 & \cdots & 0\\
\vdots & \vdots & \ddots & \vdots\\
0 & 0 & \cdots & \lambda_n
\end{pmatrix}.
\]
\end{definition}

\begin{remarque}
Les puissances d’une matrice diagonale se calculent très facilement :
\[
D^k=\diag(\lambda_1^k,\dots,\lambda_n^k).
\]
\end{remarque}

\subsection{Définition de la diagonalisation}

\begin{definition}
Un endomorphisme \(u\in\Lin(E,E)\) est dit \textbf{diagonalisable} s’il existe une base de \(E\) formée de vecteurs propres de \(u\).

Une matrice \(A\in M_n(\K)\) est dite \textbf{diagonalisable} s’il existe \(P\in GL_n(\K)\) et une matrice diagonale \(D\) telles que
\[
A=PDP^{-1}.
\]
\end{definition}

\begin{proposition}
Une matrice est diagonalisable si et seulement si elle est la matrice d’un endomorphisme diagonalisable dans une certaine base.
\end{proposition}

\begin{proof}
Dire que la matrice de \(u\) dans une base est diagonale signifie exactement que les vecteurs de cette base sont des vecteurs propres.
\end{proof}

\subsection{Critère fondamental}

\begin{theoreme}
Un endomorphisme \(u\) est diagonalisable si et seulement si
\[
E=\bigoplus_{\lambda\in\Sp(u)} E_\lambda.
\]
\end{theoreme}

\begin{proof}
Si \(u\) est diagonalisable, il existe une base de vecteurs propres. En regroupant les vecteurs associés à une même valeur propre, on obtient une base adaptée à la somme directe des sous-espaces propres.

Réciproquement, si
\[
E=\bigoplus_{\lambda\in\Sp(u)} E_\lambda,
\]
il suffit de prendre une base de chaque sous-espace propre et de les réunir. On obtient une base de \(E\) formée de vecteurs propres.
\end{proof}

\begin{corollaire}
Un endomorphisme \(u\) est diagonalisable si et seulement si la somme des dimensions des sous-espaces propres vaut \(\dim(E)\).
\end{corollaire}

\begin{proof}
Les sous-espaces propres associés à des valeurs propres distinctes sont en somme directe. Donc leur somme est directe, et elle est égale à \(E\) si et seulement si la somme de leurs dimensions vaut \(\dim(E)\).
\end{proof}

\section{Critères pratiques de diagonalisation}

\subsection{Cas des valeurs propres distinctes}

\begin{theoreme}
Si un endomorphisme de dimension \(n\) possède \(n\) valeurs propres distinctes, alors il est diagonalisable.
\end{theoreme}

\begin{proof}
Déjà vu comme conséquence du théorème d’indépendance des vecteurs propres.
\end{proof}

\subsection{Critère via les multiplicités}

\begin{theoreme}
Un endomorphisme \(u\) est diagonalisable si et seulement si :
\begin{enumerate}[label=\arabic*)]
    \item son polynôme caractéristique est scindé sur \(\K\) ;
    \item pour toute valeur propre \(\lambda\),
    \[
    \dim(E_\lambda)=m_\lambda.
    \]
\end{enumerate}
\end{theoreme}

\begin{proof}
Si \(u\) est diagonalisable, alors dans une base propre sa matrice est diagonale, donc le polynôme caractéristique est scindé et les multiplicités géométriques coïncident avec les multiplicités algébriques.

Réciproquement, si le polynôme caractéristique est scindé et si
\[
\dim(E_\lambda)=m_\lambda
\]
pour toute valeur propre, alors la somme des dimensions des sous-espaces propres vaut la somme des multiplicités algébriques, donc \(n\). Les sous-espaces propres étant en somme directe, leur somme vaut \(E\), donc \(u\) est diagonalisable.
\end{proof}

\section{Exemples de diagonalisation}

\subsection{Exemple simple}

\begin{exemple}
Soit
\[
A=
\begin{pmatrix}
2 & 0\\
0 & 3
\end{pmatrix}.
\]
Elle est déjà diagonale, donc diagonalisable.
\end{exemple}

\subsection{Exemple non diagonal mais diagonalisable}

\begin{exemple}
Soit
\[
A=
\begin{pmatrix}
1 & 1\\
0 & 2
\end{pmatrix}.
\]
Son polynôme caractéristique vaut
\[
\chiA(X)=(X-1)(X-2).
\]
Elle a deux valeurs propres distinctes, donc elle est diagonalisable.
\end{exemple}

\subsection{Exemple non diagonalisable}

\begin{exemple}
Soit
\[
A=
\begin{pmatrix}
1 & 1\\
0 & 1
\end{pmatrix}.
\]
Alors
\[
\chiA(X)=(X-1)^2.
\]
L’unique valeur propre est \(1\). Le sous-espace propre est
\[
E_1=\Ker(A-I_2)=
\Ker\begin{pmatrix}
0 & 1\\
0 & 0
\end{pmatrix}
=
\left\{
\begin{pmatrix}
x\\
0
\end{pmatrix}
: x\in\R
\right\}.
\]
Donc
\[
\dim(E_1)=1<2.
\]
La matrice n’est pas diagonalisable.
\end{exemple}

\section{Polynômes d’un endomorphisme}

\subsection{Définition}

\begin{definition}
Soit \(u\in\Lin(E,E)\) et
\[
P(X)=a_0+a_1X+\cdots+a_nX^n\in\K[X].
\]
On définit
\[
P(u)=a_0\id_E+a_1u+\cdots+a_nu^n.
\]
\end{definition}

De même, pour une matrice \(A\in M_n(\K)\),
\[
P(A)=a_0I_n+a_1A+\cdots+a_nA^n.
\]

\subsection{Valeurs propres et polynômes}

\begin{proposition}
Si \(x\) est vecteur propre de \(u\) associé à la valeur propre \(\lambda\), alors
\[
P(u)(x)=P(\lambda)x.
\]
\end{proposition}

\begin{proof}
On a
\[
u(x)=\lambda x,
\qquad
u^2(x)=\lambda^2x,
\qquad \dots,\qquad
u^n(x)=\lambda^n x.
\]
Donc
\[
P(u)(x)=a_0x+a_1\lambda x+\cdots+a_n\lambda^n x=P(\lambda)x.
\]
\end{proof}

\begin{corollaire}
Si \(P(u)=0\), alors toute valeur propre \(\lambda\) de \(u\) vérifie
\[
P(\lambda)=0.
\]
\end{corollaire}

\begin{proof}
Si \(x\neq 0\) est vecteur propre associé à \(\lambda\), alors
\[
0=P(u)(x)=P(\lambda)x.
\]
Comme \(x\neq 0\), on obtient
\[
P(\lambda)=0.
\]
\end{proof}

\subsection{Polynôme annulateur}

\begin{definition}
Un polynôme \(P\in\K[X]\) est appelé \textbf{polynôme annulateur} de \(u\) si
\[
P(u)=0.
\]
\end{definition}

\begin{remarque}
La connaissance d’un polynôme annulateur est très utile pour étudier la réduction.
\end{remarque}

\section{Critère de diagonalisation via un polynôme scindé simple}

\begin{theoreme}
Si un endomorphisme \(u\) admet un polynôme annulateur scindé à racines simples, alors \(u\) est diagonalisable.
\end{theoreme}

\begin{proof}
Supposons
\[
P(X)=(X-\lambda_1)\cdots(X-\lambda_r)
\]
avec les \(\lambda_i\) deux à deux distincts, et
\[
P(u)=0.
\]
On montre alors, par décomposition en facteurs premiers et théorème de décomposition des noyaux pour des polynômes premiers entre eux, que
\[
E=\Ker(u-\lambda_1\id_E)\oplus\cdots\oplus\Ker(u-\lambda_r\id_E).
\]
Autrement dit,
\[
E=E_{\lambda_1}\oplus\cdots\oplus E_{\lambda_r},
\]
ce qui montre que \(u\) est diagonalisable.
\end{proof}

\begin{remarque}
C’est un critère très puissant.
\end{remarque}

\section{Trigonalisabilité}

\subsection{Définition}

\begin{definition}
Un endomorphisme \(u\) est dit \textbf{trigonalisable} s’il existe une base de \(E\) dans laquelle sa matrice est triangulaire supérieure.

Une matrice \(A\in M_n(\K)\) est dite \textbf{trigonalisable} s’il existe \(P\in GL_n(\K)\) et \(T\) triangulaire supérieure telles que
\[
A=PTP^{-1}.
\]
\end{definition}

\subsection{Lien avec le polynôme caractéristique}

\begin{theoreme}
Si une matrice est trigonalisable sur \(\K\), alors son polynôme caractéristique est scindé sur \(\K\).
\end{theoreme}

\begin{proof}
Une matrice triangulaire a pour valeurs propres ses coefficients diagonaux. Donc son polynôme caractéristique est un produit de facteurs
\[
(X-\lambda_i).
\]
\end{proof}

\begin{theoreme}
Sur \(\C\), toute matrice carrée est trigonalisable.
\end{theoreme}

\begin{proof}
Le polynôme caractéristique d’une matrice complexe est scindé sur \(\C\), d’après le théorème fondamental de l’algèbre. On montre alors par récurrence sur la dimension qu’il existe une base dans laquelle la matrice est triangulaire.
\end{proof}

\begin{remarque}
Sur \(\R\), une matrice n’est pas forcément trigonalisable, car son polynôme caractéristique peut ne pas être scindé.
\end{remarque}

\section{Décomposition en sous-espaces propres}

\subsection{Somme directe des sous-espaces propres}

\begin{proposition}
Si \(\lambda_1,\dots,\lambda_r\) sont des valeurs propres deux à deux distinctes de \(u\), alors
\[
E_{\lambda_1}\oplus\cdots\oplus E_{\lambda_r}
\]
est une somme directe.
\end{proposition}

\begin{proof}
C’est une reformulation du fait que des vecteurs propres associés à des valeurs propres distinctes sont libres.
\end{proof}

\subsection{Dimension de la somme des sous-espaces propres}

\begin{corollaire}
On a
\[
\dim(E_{\lambda_1}+\cdots+E_{\lambda_r})
=
\dim(E_{\lambda_1})+\cdots+\dim(E_{\lambda_r}).
\]
\end{corollaire}

\begin{proof}
La somme est directe.
\end{proof}

\section{Cas particuliers importants}

\subsection{Endomorphisme involutif}

\begin{proposition}
Si
\[
u^2=\id_E
\]
et si \(\mathrm{char}(\K)\neq 2\), alors \(u\) est diagonalisable, avec valeurs propres dans \(\{-1,1\}\).
\end{proposition}

\begin{proof}
Le polynôme
\[
X^2-1=(X-1)(X+1)
\]
annule \(u\). Il est scindé à racines simples, donc \(u\) est diagonalisable.
\end{proof}

\subsection{Projecteur}

\begin{proposition}
Si
\[
u^2=u,
\]
alors \(u\) est diagonalisable, avec valeurs propres dans \(\{0,1\}\).
\end{proposition}

\begin{proof}
Le polynôme
\[
X^2-X=X(X-1)
\]
annule \(u\). Il est scindé à racines simples, donc \(u\) est diagonalisable.
\end{proof}

\section{Méthode pratique pour diagonaliser}

\begin{methode}
Pour étudier la diagonalisabilité d’une matrice \(A\) :
\begin{enumerate}
    \item calculer le polynôme caractéristique ;
    \item déterminer les valeurs propres ;
    \item calculer les sous-espaces propres ;
    \item comparer les dimensions des sous-espaces propres aux multiplicités algébriques ;
    \item conclure.
\end{enumerate}
\end{methode}

\begin{methode}
Si \(A\) possède \(n\) valeurs propres distinctes dans \(\K\), alors elle est immédiatement diagonalisable.
\end{methode}

\begin{methode}
Si l’on connaît un polynôme annulateur scindé à racines simples, cela suffit à conclure à la diagonalisabilité.
\end{methode}

\section{Erreurs classiques}

\begin{itemize}
    \item Confondre vecteur propre et sous-espace propre.
    \item Oublier que le vecteur nul n’est pas un vecteur propre.
    \item Croire qu’avoir une seule valeur propre empêche toujours la diagonalisabilité.
    \item Oublier de vérifier que le polynôme caractéristique est scindé sur le corps de base.
    \item Confondre multiplicité algébrique et multiplicité géométrique.
    \item Penser qu’une matrice triangulaire est toujours diagonalisable.
    \item Oublier que
    \[
    \dim(E_\lambda)\le m_\lambda.
    \]
\end{itemize}

\section{Résumé du chapitre}

\begin{itemize}
    \item Une valeur propre \(\lambda\) vérifie
    \[
    u(x)=\lambda x
    \]
    pour un vecteur non nul \(x\).
    \item Le sous-espace propre associé à \(\lambda\) est
    \[
    E_\lambda=\Ker(u-\lambda\id_E).
    \]
    \item Les valeurs propres d’une matrice sont les racines de son polynôme caractéristique.
    \item Des vecteurs propres associés à des valeurs propres distinctes sont libres.
    \item Une matrice est diagonalisable si et seulement si l’espace est somme directe de ses sous-espaces propres.
    \item Si une matrice a \(n\) valeurs propres distinctes, elle est diagonalisable.
    \item Une matrice est diagonalisable si et seulement si son polynôme caractéristique est scindé et si, pour chaque valeur propre, multiplicité algébrique et multiplicité géométrique coïncident.
    \item Si un polynôme annulateur scindé à racines simples annule \(u\), alors \(u\) est diagonalisable.
    \item Sur \(\C\), toute matrice est trigonalisable.
\end{itemize}

\section*{Exercices d’application immédiate}

\begin{enumerate}[label=\textbf{Exercice \arabic*.}]
    \item Déterminer les valeurs propres et sous-espaces propres de
    \[
    \begin{pmatrix}
    2 & 0\\
    0 & 3
    \end{pmatrix}.
    \]

    \item Étudier la diagonalisabilité de
    \[
    \begin{pmatrix}
    1 & 1\\
    0 & 2
    \end{pmatrix}.
    \]

    \item Étudier la diagonalisabilité de
    \[
    \begin{pmatrix}
    1 & 1\\
    0 & 1
    \end{pmatrix}.
    \]

    \item Montrer que des vecteurs propres associés à des valeurs propres distinctes sont libres.

    \item Montrer qu’un projecteur est diagonalisable.

    \item Montrer qu’un endomorphisme \(u\) vérifiant \(u^2=\id_E\) est diagonalisable si \(\mathrm{char}(\K)\neq 2\).

    \item Calculer le polynôme caractéristique de
    \[
    \begin{pmatrix}
    2 & 1 & 0\\
    0 & 2 & 0\\
    0 & 0 & 3
    \end{pmatrix}.
    \]

    \item Déterminer si la matrice précédente est diagonalisable.

    \item Montrer que si une matrice a \(n\) valeurs propres distinctes, alors elle est diagonalisable.

    \item Sur \(\C\), montrer qu’une matrice carrée est toujours trigonalisable.
\end{enumerate}

\end{document}